% !TeX program = xelatex
\documentclass[12pt, a4paper, titlepage]{article}
\usepackage{xeCJK}
\usepackage{amsmath, amssymb}
\usepackage{geometry}
\usepackage{longtable}
\usepackage{booktabs}
\usepackage{tabularx}
\usepackage{caption}
\usepackage{silence}
\WarningFilter{caption}{Unknown document class}
\WarningFilter{xeCJK}{Redefining CJKfamily}
\WarningFilter{relsize}{Failed to get list}

% -- Fonts --
\setmainfont{Times New Roman}
\setCJKmainfont{Hiragino Mincho ProN}

\geometry{a4paper, margin=2.5cm}

\title{高精度界面捕捉法における寄生流れ抑制のための数値指針— CCD法、Rhie-Chow補正、およびBalanced-Force条件の統合 —}
\author{}
\date{}

\begin{document}

\maketitle

\section*{1. 緒論：問題の本質}
気液二相流の数値シミュレーション、特に表面張力が支配的な微細流体現象において、\textbf{寄生流れ（Spurious Currents）}の抑制は最優先課題である。寄生流れの本質は、界面における表面張力項 $\mathbf{f}_\sigma$ と圧力勾配項 $-\nabla p$ の数値的な不整合にある。本報告では、結合コンパクト差分法（CCD法）を用いた高精度スキームにおいて、この不整合を理論的に排除するための実装指針を提示する。

\section*{2. 数値的課題と解決の論理}

\subsection*{2.1 なぜRhie-Chow補正が必要か}
コロケート格子（Collocated Grid）を採用する場合、圧力と速度のデカップリング（チェッカーボード不安定）が不可避である。\\
課題: 中心差分による圧力勾配計算は $2\Delta x$ モジュレーションを検知できず、非物理的な圧力振動を招く。\\
解決: Rhie-Chow補正により、面速度 $u_f$ に圧力勾配の差に基づく散逸項を導入する。これにより、CCD法のような高次スキームにおいても、格子点間の情報伝達の整合性を保ち、連続の式を厳密に満足させることが可能となる。

\subsection*{2.2 なぜCCD法で曲率を求めるだけでは不十分か}
高精度な曲率 $\kappa$ の算出は、表面張力の「入力精度」を上げるが、それだけでは「力のバランス」は保たれない。\\
課題: 非常に正確な曲率 $\kappa$ を得ても、その勾配 $\nabla \psi$ を計算する演算子と、圧力勾配 $\nabla p$ を計算する演算子が異なれば、静止状態で両者がキャンセルされず、誤差が速度として現れる。\\
解決: Balanced-Force（バランスド・フォース）条件の導入。

\section*{3. Balanced-Force 離散化の提案}
寄生流れを機械精度レベルまで抑制するため、以下の実装スキームを提案する。

\begin{itemize}
    \item \textbf{演算子の統一}: 運動量方程式における $\nabla p$ と、表面張力項内の $\nabla \psi$ の計算に、全く同一のCCD演算子を適用する。
    \item \textbf{曲率の高次近似}: 界面関数 $\psi$ から法線ベクトル $\mathbf{n}$、曲率 $\kappa$ を求める全プロセスをCCD法（$O(h^6)$）で行い、滑らかな分布を得る。
    \item \textbf{表面張力の定式化}: CSFモデルをベースとし、以下の形式で実装する。
    \begin{equation*}
        \mathbf{f}_\sigma = \sigma \kappa [ \nabla^{CCD} \psi ]
    \end{equation*}
    ここで $[ \nabla^{CCD} ]$ は圧力勾配計算と同一のサブルーチンを指す。
\end{itemize}

\section*{4. 今後の実装ステップ（Action Plan）}

\subsection*{Phase 1: 基盤の実装}
\begin{itemize}
    \item [ ] CCD演算子の共通化: 既存の圧力勾配計算ルーチンを、表面張力の勾配計算にも適用できるようモジュール化する。
    \item [ ] Rhie-Chow項の再定義: CCD法による圧力勾配を考慮した新しい面速度補間ルーチンの実装。
\end{itemize}

\subsection*{Phase 2: 検証（Benchmark）}
\begin{itemize}
    \item [ ] 静止液滴テスト（Static Drop）: 寄生流れの大きさが時間とともに増大しないか、マシンスケール（$10^{-12}$以下）に収まっているかを確認。
    \item [ ] Laplace圧の検証: 数値的な圧力跳び $\Delta p$ が理論値 $\sigma/R$ と一致するかを確認。
\end{itemize}

\subsection*{Phase 3: 応用}
\begin{itemize}
    \item [ ] 液滴の衝突・合体シミュレーション: 高精度界面記述がトポロジー変化に与える影響の評価。
\end{itemize}

\section*{5. 論文にまとめるべき論点}
\begin{itemize}
    \item \textbf{CCD法とBalanced-Forceの親和性}: 高次差分を用いることで、低次スキームよりも寄生流れ抑制が困難になる（誤差が鋭敏に現れる）性質を逆手に取り、演算子の一致がいかに劇的な効果を生むかを定量的に示す。
    \item \textbf{誤差評価}: 界面の厚み $\epsilon$ に対する収束性と、格子解像度 $h$ に対する精度の関係を整理。
    \item \textbf{計算効率}: Rhie-Chow補正を伴うCCD法が、反復解法の収束性に与えるポジティブな影響の報告。
\end{itemize}

\section*{結論}
Rhie-Chow補正は「格子の幾何学的欠陥」を補い、Balanced-Forceは「物理項間の数値的不整合」を解消する。両者をCCD法という強力な枠組みで統合することで、極めて信頼性の高い二相流ソルバーが実現される。

\end{document}